\chapter{Obstacle problems and variational inequalities}\label{sec:obstacle-problems}
This chapter begins the obstacle-problem branch with a model variational formulation. The Sobolev and weak-solution background is developed in \cite{SimonPDE}, Chapters 5--9. A dedicated source corpus for obstacle regularity has not yet been added; the statements below are elementary calculations, not a survey of that theory.

\begin{definition}[The classical scalar obstacle problem]\label{def:obstacle-problem}
Let $U\subset\mathbb R^d$ be bounded with Lipschitz boundary, $f\in L^2(U)$, and $g,\psi\in H^1(U)$. Assume the convex admissible set
\[
 K=\{v\in H^1(U):v-g\in H^1_0(U),\ v\geq\psi\text{ almost everywhere}\}
\]
is nonempty. The obstacle problem is to minimize
$\mathcal E(v)=\frac12\int_U|Dv|^2-\int_Ufv$
over $K$. The inequality is a unilateral constraint. This entry defines the model; existence and regularity are separate questions.
\end{definition}

\begin{proposition}[Variational inequality for the obstacle problem]\label{prop:obstacle-inequality}
For $u\in K$, minimization of $\mathcal E$ over $K$ is equivalent to
\[
 \int_U Du\cdot D(v-u)\geq\int_U f(v-u)\qquad\text{for every }v\in K.
\]
\end{proposition}
\begin{proof}
If $u$ minimizes and $v\in K$, convexity gives $u+t(v-u)\in K$ for $0\leq t\leq1$. Its energy has nonnegative right derivative at zero, which is exactly the inequality. Conversely, expand the quadratic energy:
\[
 \mathcal E(v)-\mathcal E(u)
 =\int_U Du\cdot D(v-u)-\int_U f(v-u)+\tfrac12\int_U|D(v-u)|^2.
\]
Every term on the right after combining the first two is nonnegative under the stated inequality. Hence $u$ minimizes.
\end{proof}

\begin{remark}[Planned free-boundary developments]\label{rem:obstacle-program}
When $u$ and $\psi$ have continuous representatives, the contact set is $\{u=\psi\}$ and the interior free boundary is $\partial\{u>\psi\}\cap U$. Future sections will develop existence, optimal regularity, blow-ups, nondegeneracy, and the regular and singular parts of the free boundary. Thin-obstacle problems and geometric obstacles require distinct competitor spaces and hypotheses. No free-boundary regularity theorem is claimed at this stage.
\end{remark}
